\begin{conclusion}[微分表示对属 \texorpdfstring{$49$}{49} 的 \texorpdfstring{$S_4$}{S4}--Hurwitz 链全部商曲线亏格的统一锁定]\label{con:xi-time-part9n-s4-all-subgroup-genera}
设
\[
V:=H^0(X,\Omega_X^1)
\cong 5\cdot \mathrm{sgn}\oplus 4\cdot V_2\oplus 3\cdot V_3\oplus 9\cdot V_3'
\]
为属 \(49\) 的 \texorpdfstring{$S_4$}{S4}--Hurwitz 曲线 \(X\) 的微分表示。记
\[
H:=X/A_4,\qquad
C_6:=X/V_4,\qquad
Z:=X/D_8,\qquad
C_F:=X/S_3,\qquad
Y:=X/C_4,
\]
并令 \(C_2=\langle\tau\rangle\)、\(C_3=\langle\sigma_3\rangle\)、\(C_4=\langle\sigma_4\rangle\) 分别表示由换位、\(3\)-循环与 \(4\)-循环生成的循环子群。则全部中间商曲线的亏格满足
\[
g(X/S_4)=0,\qquad
g(H)=g(X/A_4)=5,\qquad
g(C_6)=g(X/V_4)=13,\qquad
g(Z)=g(X/D_8)=4,
\]
\[
g(C_F)=g(X/S_3)=3,\qquad
g(Y)=g(X/C_4)=13,\qquad
g(X/C_3)=17,\qquad
g(X/C_2)=19.
\]
特别地，\(C_3\)-商与换位商的亏格无需额外 Hurwitz 计算，已由同一微分表示数据唯一锁定。
\end{conclusion}
\begin{proof}
对任意子群 \(K\le S_4\)，商曲线亏格满足
\[
g(X/K)=\dim V^K.
\]
记
\[
\rho_K:=\operatorname{Ind}_{K}^{S_4}\mathbf 1.
\]
由 Frobenius 互反，
\[
\dim(\pi^K)=\langle \pi,\rho_K\rangle_{S_4}
\]
对每个不可约表示 \(\pi\) 都成立。另一方面，定理 \ref{thm:killo-s4-burnside-kani-rosen-prym-square} 给出
\[
\rho_{A_4}=\mathbf 1\oplus \mathrm{sgn},\qquad
\rho_{V_4}=\mathbf 1\oplus \mathrm{sgn}\oplus 2V_2,\qquad
\rho_{D_8}=\mathbf 1\oplus V_2,
\]
\[
\rho_{S_3}=\mathbf 1\oplus V_3,\qquad
\rho_{C_4}=\mathbf 1\oplus V_2\oplus V_3',\qquad
\rho_{C_2}=\mathbf 1\oplus V_2\oplus 2V_3\oplus V_3',
\]
\[
\rho_{C_3}=\mathbf 1\oplus \mathrm{sgn}\oplus V_3\oplus V_3',\qquad
\rho_{S_4}=\mathbf 1.
\]
将这些分解与 \(V\) 的 Chevalley--Weil 分解逐项配对，便得到
\[
g(X/A_4)=5,\quad
g(X/V_4)=5+2\cdot4=13,\quad
g(X/D_8)=4,\quad
g(X/S_3)=3,
\]
\[
g(X/C_4)=4+9=13,\quad
g(X/C_3)=5+3+9=17,\quad
g(X/C_2)=4+2\cdot3+9=19.
\]
又因 \(V\) 中不含平凡表示，故
\[
g(X/S_4)=\dim V^{S_4}=0.
\]
证毕。
\end{proof}

\begin{conclusion}[属 \texorpdfstring{$49$}{49} 的 \texorpdfstring{$S_4$}{S4}--Hurwitz 链中间商映射总分歧度的 \texorpdfstring{$24$}{24} 倍数谱]\label{con:xi-time-part9n-s4-total-ramification-spectrum}
对上述各子群 \(K\le S_4\)，记
\[
\pi_K:X\to X/K
\]
为自然商映射，并令
\[
R_K:=\sum_{P\in X}\bigl(e_P(\pi_K)-1\bigr)
\]
为其总分歧度。则有显式公式
\[
R_{S_4}=144,\qquad
R_{A_4}=0,\qquad
R_{V_4}=0,\qquad
R_{D_8}=48,
\]
\[
R_{S_3}=72,\qquad
R_{C_4}=0,\qquad
R_{C_3}=0,\qquad
R_{C_2}=24.
\]
因此全部非零总分歧度恰落在
\[
24\cdot\{1,2,3,6\},
\]
并形成严格链
\[
24<48<72<144.
\]
\end{conclusion}
\begin{proof}
由上一定理与 Chevalley--Weil 分解可知 \(g(X)=49\)，故
\[
2g(X)-2=96.
\]
对任意 \(K\le S_4\)，Riemann--Hurwitz 公式给出
\[
2g(X)-2=|K|\bigl(2g(X/K)-2\bigr)+R_K,
\]
即
\[
R_K=96-|K|\bigl(2g(X/K)-2\bigr).
\]
将结论 \ref{con:xi-time-part9n-s4-all-subgroup-genera} 的亏格值与
\[
|S_4|=24,\quad |A_4|=12,\quad |V_4|=4,\quad |D_8|=8,\quad
|S_3|=6,\quad |C_4|=4,\quad |C_3|=3,\quad |C_2|=2
\]
逐一代入，立即得到
\[
R_{S_4}=144,\quad R_{A_4}=0,\quad R_{V_4}=0,\quad R_{D_8}=48,
\]
\[
R_{S_3}=72,\quad R_{C_4}=0,\quad R_{C_3}=0,\quad R_{C_2}=24.
\]
其中 \(R_{A_4}=R_{V_4}=R_{C_3}=0\) 与换位惯性对子群 \(A_4,V_4,C_3\) 的平凡交相一致，而 \(R_{C_4}=0\) 则对应于 \(C_4\) 中不含任何换位惯性元。证毕。
\end{proof}

\begin{conclusion}[属 \texorpdfstring{$49$}{49} 的 \texorpdfstring{$S_4$}{S4}--Hurwitz Jacobian 的 \texorpdfstring{$13+9+9+9+9$}{13+9+9+9+9} 型块压缩]\label{con:xi-time-part9n-jacobian-c6-compression}
设
\[
J(X)\sim_{\QQ}J(H)\times J(Z)^2\times J(C_F)^3\times P^3,
\qquad
P=\operatorname{Prym}(Y/Z),
\]
并沿上文记号令 \(C_6=X/V_4\)。则有更压缩的同源分解
\[
\boxed{\;
J(X)\sim_{\QQ} J(C_6)\times J(C_F)^3\times P^3.
\;}
\]
等价地，原来的
\[
5+2\cdot 4+3\cdot 3+3\cdot 9=49
\]
维分解可压缩为
\[
13+9+9+9+9=49
\]
的块结构，其中 \(13\) 维块 \(J(C_6)\) 已吸收 \(J(H)\) 与 \(J(Z)^2\) 的全部几何信息。
\end{conclusion}
\begin{proof}
定理 \ref{thm:killo-s4-burnside-kani-rosen-prym-square} 给出
\[
J(C_6)\sim_{\QQ}J(H)\times J(Z)^2,
\]
其几何内容正是无分歧循环三重覆盖 \(C_6\to H\) 的 Prym 分裂
\[
\operatorname{Prym}(C_6/H)\sim_{\QQ}J(Z)^2.
\]
另一方面，定理 \ref{thm:killo-s4-genus49-jacobian-computable-isogeny} 给出
\[
J(X)\sim_{\QQ}J(H)\times J(Z)^2\times J(C_F)^3\times P^3.
\]
将前一式代入后一式，立即得到
\[
J(X)\sim_{\QQ}J(C_6)\times J(C_F)^3\times P^3.
\]
又由
\[
g(C_6)=13,\qquad 3g(C_F)=9,\qquad \dim P^3=27=9+9+9
\]
可知该分解正对应于 \(13+9+9+9+9\) 的块压缩。证毕。
\end{proof}

\endinput
